\textbf{Introduction générale}

L'intelligence artificielle générative, dont l'essor a marqué ces dernières années avec une vigueur remarquable, a profondément reconfiguré la manière dont les organisations envisagent l'accès au savoir, l'assistance aux utilisateurs et l'automatisation de leurs processus internes. Parmi les approches qui ont retenu l'attention des chercheurs comme des praticiens, les systèmes conversationnels fondés sur le paradigme \textbf{Retrieval-Augmented Generation} (RAG) se distinguent par leur capacité à pallier les insuffisances des modèles génératifs purs : en articulant génération automatique de réponses et exploitation raisonnée de bases documentaires spécialisées, ils offrent une pertinence contextuelle que les modèles seuls peinent à atteindre.

Or, dans la réalité des environnements professionnels, les besoins excèdent largement celui d'un assistant conversationnel généraliste. Les organisations sont fréquemment confrontées à la nécessité de disposer de plusieurs assistants spécialisés, chacun ancré dans un domaine précis — ressources humaines, support informatique, conformité réglementaire, documentation interne, ou encore accompagnement universitaire. Si chacun de ces assistants venait à être développé de manière isolée, les conséquences seraient inévitables : duplication des composants, disparité des mécanismes de sécurité, alourdissement de la maintenance et résistance à toute tentative d'industrialisation. C'est pourquoi la mise en place d'une plateforme centralisée — capable d'héberger plusieurs chatbots métiers tout en assurant la gouvernance des accès, la protection des données et la traçabilité de chaque traitement — représente un enjeu qui est simultanément technique, organisationnel et stratégique.

C'est dans cette perspective que prend place le présent projet, intitulé \textbf{SecureRAG Hub}. Il s'agit de concevoir et de réaliser une plateforme unifiée de chatbots métiers sécurisés, reposant sur une architecture microservices, une chaîne DevSecOps rigoureuse et des mécanismes de sécurité pilotés par l'intelligence artificielle. L'ambition est de proposer un environnement à la fois extensible et reproductible, au sein duquel plusieurs assistants spécialisés peuvent coexister, chacun doté de son propre domaine métier, de sa base documentaire dédiée, de ses règles d'accès et de ses dispositifs intelligents de contrôle. Dans le cadre de cette étude, deux chatbots servent de preuve de concept : un assistant dédié aux ressources humaines et un assistant orienté support informatique.

La problématique centrale peut se formuler en ces termes : \textbf{comment concevoir une plateforme centralisée capable d'héberger plusieurs chatbots métiers sécurisés, articulant architecture microservices, orchestration RAG, gouvernance des accès et mécanismes de sécurité pilotés par l'IA, tout en demeurant extensible, traçable et déployable au sein d'une chaîne DevSecOps ?} Cette question met en lumière l'exigence d'intégrer, dans une solution cohérente et unifiée, des dimensions que l'on traite trop souvent séparément : architecture logicielle distribuée, recherche vectorielle, génération assistée par des modèles de langage de grande taille, sécurité applicative, protection contre les attaques propres aux systèmes RAG — notamment les injections de prompts et les tentatives d'exfiltration — ainsi qu'automatisation des processus qualité et de déploiement.

L'originalité du projet tient à la convergence de plusieurs axes technologiques qui se renforcent mutuellement. L'architecture microservices structure la plateforme autour de composants clairement délimités : une API Gateway, un service d'authentification et de gestion des rôles, un gestionnaire de chatbots, un orchestrateur LLM, un composant de sécurité intelligent et un \textbf{Knowledge Hub} fondé sur Qdrant. L'approche DevSecOps, quant à elle, garantit l'intégration continue des contrôles de qualité et de sécurité à travers une chaîne d'outils comprenant GitHub Actions, Semgrep, Gitleaks, Trivy, Syft et Cosign. Enfin, la dimension \textbf{AI-Sec} introduit une sécurité véritablement pilotée par l'IA, grâce à la combinaison de règles NLP, d'un scoring sémantique par embeddings et d'une évaluation par \textbf{LLM-as-judge} — formant ce que l'on nomme dans ce travail le paradigme AI-Sec à trois niveaux.

Le choix de ce sujet se justifie pleinement dans le cadre d'un master orienté systèmes d'information et développement intelligent, tant il mobilise des compétences transversales qui traversent l'architecture logicielle, l'intelligence artificielle, la cybersécurité, le déploiement conteneurisé et l'automatisation des chaînes de livraison. Mais au-delà de son intérêt académique, ce projet répond à des préoccupations très concrètes du monde professionnel : sécurisation des données sensibles, maîtrise fine des accès selon les rôles, robustesse face aux attaques visant les systèmes conversationnels, et capacité à reproduire fidèlement un déploiement dans un environnement contrôlé.

L'objectif général de ce mémoire est donc de présenter la conception, la mise en œuvre et l'évaluation de cette plateforme unifiée. Plus précisément, il convient de définir une architecture technique adaptée, de constituer une base documentaire vectorielle mutualisée, d'orchestrer les échanges entre utilisateurs, chatbots et modèles de langage, de développer un composant de sécurité intelligent, puis d'intégrer l'ensemble dans une chaîne DevSecOps aboutissant à un déploiement sur Kubernetes local. La faisabilité et la pertinence de cette approche sont vérifiées à travers la démonstration concrète que constituent les deux chatbots retenus.

Le présent mémoire s'articule en cinq chapitres. Le premier expose le contexte général, la problématique, les objectifs et le périmètre du projet. Le deuxième est consacré à la conception de la solution, en portant l'accent sur l'architecture microservices, les composants fonctionnels et les flux de traitement. Le troisième décrit les choix technologiques retenus ainsi que la mise en œuvre concrète des différents services de la plateforme. Le quatrième détaille l'intégration des mécanismes de sécurité et de la chaîne DevSecOps. Le cinquième, enfin, présente le déploiement, l'expérimentation et l'évaluation de la solution. Une conclusion générale clôt le mémoire en récapitulant les résultats obtenus, les limites observées et les perspectives qui se dessinent.

À travers ce travail, nous cherchons ainsi à démontrer qu'il est possible de concevoir une plateforme conversationnelle moderne, modulaire et rigoureusement sécurisée, réconciliant pertinence métier, gouvernance des accès, sécurité pilotée par l'IA et exigences d'industrialisation logicielle — non pas comme un assemblage de contraintes contradictoires, mais comme une architecture pensée dans sa cohérence dès l'origine.
